% GENERATED FILE. Edit canonical lessons, then run npm run generate:books.
% canonical-source-hash: fnv1a64:cf9f564d0b42733d
% canonical-lessons: ZH-C243-qihou, ZH-C243-xuehua, ZH-C243-bianhua, ZH-C243-jiebing, ZH-C243-wendu

\chapter{Changing Seasons}
\label{ch:zh-changing-seasons}
\hlchaptermodality{243}

\begin{chapteropening}
\textbf{By the end of this chapter:} \emph{I can describe the climate, snowflakes, the weather changing, ice on the lake and the temperature.}

Complete the last lesson of chapter 243: I can describe the climate, snowflakes, the weather changing, ice on the lake and the temperature.
\end{chapteropening}

\section[\texorpdfstring{\zh{气候} (qìhòu)}{气候 (qìhòu)}]{\zh{气候} (qìhòu) — climate}

\label{lesson:ZH-C243-qihou}



\begin{quote}
Before the new one: say the Mandarin for a bug, then the Mandarin for a mosquito.
\end{quote}

\subsection*{You'll want to know: \zh{气候}}
\textbf{\zh{气候}} — \emph{qìhòu} — \textquotedblleft{}climate\textquotedblright{}.

\begin{quote}
\textbf{\zh{西边的气候不好风很大}} \emph{xībiān de qìhòu bù hǎo fēng hěn dà} — The climate in the west isn't good: it's very windy.
\end{quote}

\begin{cousinweb}
\textbf{\zh{气候}} (\emph{qìhòu}) is built from \zh{气} 'air' and \zh{候} 'a season', both inherited from Middle Chinese, \emph{khj+jH} and \emph{huwH} (Guangyun, in Baxter's transcription), which glosses the second 'to wait on'; the sense of a season is not given there.
\end{cousinweb}

\begin{grammarlens}[title={What you've built}]
One more word for describing things.
\end{grammarlens}

\subsection*{Your turn}
\begin{itemize}
\raggedright
  \item \emph{Say it:} \emph{qìhòu}
  \item \emph{Say it:} \emph{qìhòu}, once more
  \item \emph{Say it:} \emph{xībiān de qìhòu bù hǎo fēng hěn dà} — The climate in the west isn't good: it's very windy.
  \item \emph{Recall:} say the Mandarin for a bug, then the Mandarin for a mosquito, then say \emph{qìhòu} again
\end{itemize}

\begin{tcolorbox}[breakable,colback=teal!4,colframe=teal!35!black,title={Before you move on}]
What does \zh{气候} mean? (\textquotedblleft{}Climate\textquotedblright{}.) Say it once more.
\end{tcolorbox}
\section[\texorpdfstring{\zh{雪花} (xuěhuā)}{雪花 (xuěhuā)}]{\zh{雪花} (xuěhuā) — snowflake}

\label{lesson:ZH-C243-xuehua}



\begin{quote}
Before the new one: say the Mandarin for a mosquito, then the Mandarin for climate.
\end{quote}

\subsection*{You'll want to know: \zh{雪花}}
\textbf{\zh{雪花}} — \emph{xuěhuā} — \textquotedblleft{}snowflake\textquotedblright{}.

\begin{quote}
\textbf{\zh{今天多云外边有雪花}} \emph{jīntiān duōyún wàibian yǒu xuěhuā} — It's cloudy today, and there are snowflakes outside.
\end{quote}

\begin{cousinweb}
\textbf{\zh{雪花}} (\emph{xuěhuā}) is built from \zh{雪} 'snow' and \zh{花} 'flower', both inherited from Middle Chinese, \emph{sjwet} and \emph{xwae} (Guangyun, in Baxter's transcription).
\end{cousinweb}

\begin{grammarlens}[title={What you've built}]
One more word for describing things.
\end{grammarlens}

\subsection*{Your turn}
\begin{itemize}
\raggedright
  \item \emph{Say it:} \emph{xuěhuā}
  \item \emph{Say it:} \emph{xuěhuā}, once more
  \item \emph{Say it:} \emph{jīntiān duōyún wàibian yǒu xuěhuā} — It's cloudy today, and there are snowflakes outside.
  \item \emph{Recall:} say the Mandarin for a mosquito, then the Mandarin for climate, then say \emph{xuěhuā} again
\end{itemize}

\begin{tcolorbox}[breakable,colback=teal!4,colframe=teal!35!black,title={Before you move on}]
What does \zh{雪花} mean? (\textquotedblleft{}Snowflake\textquotedblright{}.) Say it once more.
\end{tcolorbox}
\section[biànhuà]{biànhuà — to change over time; a change}

\label{lesson:ZH-C243-bianhua}



\begin{quote}
Before the new one: say the Mandarin for climate, then the Mandarin for a snowflake.
\end{quote}

\subsection*{You'll want to know: biànhuà}
\emph{biànhuà} — \textquotedblleft{}to change over time; a change\textquotedblright{}.

\begin{quote}
\emph{zhèr de tiānqì biànhuà hěn dà} — The weather here changes a lot.
\end{quote}

Its characters come later; for now it is a word for the ear.

\begin{cousinweb}
\emph{biànhuà} is built from \emph{biàn} 'to alter' and \emph{huà} 'to transform', both inherited from Middle Chinese, \emph{pjenH} and \emph{xwaeH} (Guangyun, in Baxter's transcription), which glosses each with the other.
\end{cousinweb}

\begin{grammarlens}[title={What you've built}]
One more word for describing things.
\end{grammarlens}

\subsection*{Your turn}
\begin{itemize}
\raggedright
  \item \emph{Say it:} \emph{biànhuà}
  \item \emph{Say it:} \emph{biànhuà}, once more
  \item \emph{Say it:} \emph{zhèr de tiānqì biànhuà hěn dà} — The weather here changes a lot.
  \item \emph{Recall:} say the Mandarin for climate, then the Mandarin for a snowflake, then say \emph{biànhuà} again
\end{itemize}

\begin{tcolorbox}[breakable,colback=teal!4,colframe=teal!35!black,title={Before you move on}]
What does biànhuà mean? (\textquotedblleft{}To change over time; a change\textquotedblright{}.) Say it once more.
\end{tcolorbox}
\section[jiébīng]{jiébīng — to freeze over}

\label{lesson:ZH-C243-jiebing}



\begin{quote}
Before the new one: say the Mandarin for a snowflake, then the Mandarin for to change over time.
\end{quote}

\subsection*{You'll want to know: jiébīng}
\emph{jiébīng} — \textquotedblleft{}to freeze over\textquotedblright{}.

\begin{quote}
\emph{dōngtiān hú lǐ de shuǐ jiébīng le} — In winter the water in the lake freezes over.
\end{quote}

Its characters come later; for now it is a word for the ear.

\begin{cousinweb}
\emph{jiébīng} is built from \emph{jié} 'to knot' and \emph{bīng} 'ice', both inherited from Middle Chinese, \emph{ket} and \emph{ping} (Guangyun, in Baxter's transcription), which glosses the first 'to tie' and the second 'water frozen'.
\end{cousinweb}

\begin{grammarlens}[title={What you've built}]
One more word for describing things.
\end{grammarlens}

\subsection*{Your turn}
\begin{itemize}
\raggedright
  \item \emph{Say it:} \emph{jiébīng}
  \item \emph{Say it:} \emph{jiébīng}, once more
  \item \emph{Say it:} \emph{dōngtiān hú lǐ de shuǐ jiébīng le} — In winter the water in the lake freezes over.
  \item \emph{Recall:} say the Mandarin for a snowflake, then the Mandarin for to change over time, then say \emph{jiébīng} again
\end{itemize}

\begin{tcolorbox}[breakable,colback=teal!4,colframe=teal!35!black,title={Before you move on}]
What does jiébīng mean? (\textquotedblleft{}To freeze over\textquotedblright{}.) Say it once more.
\end{tcolorbox}
\section[wēndù]{wēndù — temperature}

\label{lesson:ZH-C243-wendu}



\begin{quote}
Before the new one: say the Mandarin for to change over time, then the Mandarin for to freeze over.
\end{quote}

\subsection*{You'll want to know: wēndù}
\emph{wēndù} — \textquotedblleft{}temperature\textquotedblright{}.

\begin{quote}
\emph{jīntiān wēndù bù gāo} — It's not very warm today.
\end{quote}

Its characters come later; for now it is a word for the ear.

\begin{cousinweb}
\emph{wēndù} is built from \emph{wēn} 'warm' and \emph{dù} 'a measure', both inherited from Middle Chinese, \emph{'won} and \emph{duH} (Guangyun, in Baxter's transcription), which glosses the second 'a rule'; the sense of a degree on a scale is not given there.
\end{cousinweb}

\begin{grammarlens}[title={What you've built}]
One more word for describing things.
\end{grammarlens}

\subsection*{Your turn}
\begin{itemize}
\raggedright
  \item \emph{Say it:} \emph{wēndù}
  \item \emph{Say it:} \emph{wēndù}, once more
  \item \emph{Say it:} \emph{jīntiān wēndù bù gāo} — It's not very warm today.
  \item \emph{Recall:} say the Mandarin for to change over time, then the Mandarin for to freeze over, then say \emph{wēndù} again
\end{itemize}

\begin{tcolorbox}[breakable,colback=teal!4,colframe=teal!35!black,title={Before you move on}]
What does wēndù mean? (\textquotedblleft{}Temperature\textquotedblright{}.) Say it once more.
\end{tcolorbox}
